% ============================================================
%  Supplementary File: Complete 24-Item Reporting Checklist
% ============================================================

\section*{Supplementary File: Complete 24-Item Reporting Checklist for Fairness in Histopathology AI}
\label{app:checklist}

This supplementary file provides the complete 24-item reporting checklist referenced in Section~\ref{sec:discussion:agenda:checklist} of the main text. The checklist is organised into five blocks (A--E) and is designed to complement---not replace---CLAIM, TRIPOD+AI, and STARD-AI guidelines.

\begin{table}[htbp]
  \centering
  \caption{Complete 24-item reporting checklist. Items marked $\dagger$ are
           \textbf{mandatory} for all submissions; items marked $\ddagger$ are
           \textbf{conditional} (required if applicable to the study).
           Block~D operationalises the six-item minimum metric floor proposed
           in this review.}
  \label{tab:complete_checklist}
  \small
  \begin{tabularx}{\linewidth}{@{}c p{10.5cm} X@{}}
    \toprule
    \textbf{Item} & \textbf{Requirement} & \textbf{Section} \\
    \midrule
    \multicolumn{3}{l}{\textbf{Block A: Tissue Provenance (5 items)}} \\
    A1$^\dagger$ & List all tissue sources (hospitals, biobanks, public datasets) with geographic location (city, country). & \ref{app:checklist:blockA} \\
    A2$^\dagger$ & Specify scanner models, manufacturers, and magnifications used for each site. & \ref{app:checklist:blockA} \\
    A3$^\dagger$ & Describe staining protocols (H\&E, IHC markers, antibody clones, incubation times) per site. & \ref{app:checklist:blockA} \\
    A4$^\ddagger$ & If multiple staining protocols were used, provide rationale and distribution across cohorts. & \ref{app:checklist:blockA} \\
    A5$^\ddagger$ & Report quality-control logs: percentage of slides excluded for artefacts, folding, or poor staining per site. & \ref{app:checklist:blockA} \\
    \midrule
    \multicolumn{3}{l}{\textbf{Block B: Pre-processing Audit (4 items)}} \\
    B1$^\dagger$ & Specify normalisation method (Reinhard, Macenko, Vahadane, etc.) with template slide or reference distribution. & \ref{app:checklist:blockB} \\
    B2$^\dagger$ & Report augmentation pipeline (types, magnitudes, probabilities) applied during training. & \ref{app:checklist:blockB} \\
    B3$^\ddagger$ & If stain normalisation was \textit{not} applied, justify the decision. & \ref{app:checklist:blockB} \\
    B4$^\ddagger$ & For foundation-model fine-tuning, specify whether encoder weights were frozen or updated, and learning rates for each. & \ref{app:checklist:blockB} \\
    \midrule
    \multicolumn{3}{l}{\textbf{Block C: Cohort Demographics (5 items)}} \\
    C1$^\dagger$ & Report per-subgroup sample sizes (patients, slides, tiles) for all protected attributes analyzed. & \ref{app:checklist:blockC} \\
    C2$^\dagger$ & Specify attribute source: self-reported, clinician-assigned, or algorithmically inferred. & \ref{app:checklist:blockC} \\
    C3$^\dagger$ & Report geographic composition: percentage of data from each continent or income level (World Bank classification). & \ref{app:checklist:blockC} \\
    C4$^\dagger$ & Disclose missingness: percentage of samples with unknown/unreported values for each demographic field. & \ref{app:checklist:blockC} \\
    C5$^\ddagger$ & For studies using public datasets, cite the specific demographic metadata audit source (e.g., \citet{norori2021addressing}) or provide extracted demographics. & \ref{app:checklist:blockC} \\
    \midrule
    \multicolumn{3}{l}{\textbf{Block D: Fairness Evaluation (6 items) --- MINIMUM REPORTING FLOOR}} \\
    D1$^\dagger$ & Report aggregate AUROC (or primary metric) separately for in-distribution and out-of-distribution test splits. & \ref{app:checklist:blockD} \\
    D2$^\dagger$ & Report at least one subgroup disparity metric: subgroup AUC gap, TPR/FPR disparity, or equalized odds difference. & \ref{app:checklist:blockD} \\
    D3$^\dagger$ & Report worst-group accuracy or AUC (minimum performance across all predefined subgroups). & \ref{app:checklist:blockD} \\
    D4$^\dagger$ & Report Expected Calibration Error (ECE) overall and, if sample size permits, stratified by subgroup. & \ref{app:checklist:blockD} \\
    D5$^\dagger$ & Conduct and report preserved-site cross-validation or equivalent site-leakage probe (e.g., site-classification accuracy on features). & \ref{app:checklist:blockD} \\
    D6$^\ddagger$ & For MIL models, report attention-level fairness: correlation between attention weights and diagnostic regions across subgroups. & \ref{app:checklist:blockD} \\
    \midrule
    \multicolumn{3}{l}{\textbf{Block E: Reproducibility (4 items)}} \\
    E1$^\dagger$ & Provide code availability status: public URL, upon request, or unavailable (with justification). & \ref{app:checklist:blockE} \\
    E2$^\ddagger$ & If code is public, include model weights, trained checkpoints, and evaluation harness. & \ref{app:checklist:blockE} \\
    E3$^\dagger$ & Specify dataset access tiers: open access, controlled access (DU required), or unavailable (with citation for access procedure). & \ref{app:checklist:blockE} \\
    E4$^\dagger$ & Document evaluation harness: hyperparameters, random seeds, hardware specifications, and runtime for all experiments. & \ref{app:checklist:blockE} \\
    \bottomrule
  \end{tabularx}
\end{table}

\subsection*{Implementation Guidance}

\paragraph{Block A: Tissue Provenance}
Multi-site studies should provide a table listing each site with patient counts, scanner details, and staining protocols. For public datasets, cite the dataset accession number and version. Site-level metadata is critical for detecting institutional bias and assessing generalisability.

\paragraph{Block B: Pre-processing Audit}
Stain normalisation method and template selection substantially affect downstream fairness~\citep{voon2023evaluating, lin2025stain}. Authors should specify whether the template was a single reference slide, average distribution, or learned target. Augmentation pipelines should be fully documented, as colour augmentation can partially substitute for normalisation~\citep{franchet2024bias}.

\paragraph{Block C: Cohort Demographics}
Demographic reporting should follow the principle of ``maximum feasible granularity'': race/ethnicity categories should match the source instrument (e.g., US Census categories for US cohorts), with an ``unknown'' category for missing data. Geographic composition should use World Bank income classifications (high-income, upper-middle-income, lower-middle-income, low-income) to enable cross-study comparison of representativeness.

\paragraph{Block D: Fairness Evaluation}
This block constitutes the \textbf{minimum reporting floor} for any study claiming to address fairness. All six items are mandatory for peer-reviewed publication. The subgroup calibration gap (D4) may be omitted if subgroup sample sizes are $<$50, but this limitation must be explicitly disclosed. For MIL models, attention-level fairness (D6) is emerging best practice but not yet universally feasible.

\paragraph{Block E: Reproducibility}
Code and data availability should follow FAIR principles (Findable, Accessible, Interoperable, Reusable). For datasets with access restrictions, authors should provide a ``data access statement'' with contact information and typical turnaround time. Evaluation harness documentation is critical for cross-study comparison: hyperparameters, seeds, and hardware can substantially affect reported fairness metrics~\citep{zong2023medfair}.

\subsection*{Editorial Implementation}

We recommend the following implementation pathway:

\begin{enumerate}
  \item \textbf{Journal submission}: Require completion of Blocks A, C, D, and E as a precondition for peer review. Block B (pre-processing) may be moved to supplementary materials for word-count purposes.
  
  \item \textbf{Peer review}: Assign a dedicated ``fairness reviewer'' to verify Block D (Fairness Evaluation) completeness. Incomplete Block D should result in major revision.
  
  \item \textbf{Post-publication}: Journals should publish completed checklists as supplementary files, enabling meta-research on reporting compliance.
  
  \item \textbf{Enforcement}: After a 12-month grace period, journals should desk-reject submissions with incomplete Block D reporting.
\end{enumerate}

This checklist is released under a CC-BY 4.0 license and may be adapted by journals, conferences, and funding agencies without permission, provided the original source is cited.
